\section{Segment Data Blocks}\label{sec:segment-data-blocks}

Segment data blocks provide the information required to calculate one trajectory
segment. This includes numerical-integration parameters, vehicle configuration,
guidance rules, segment final conditions, and what to do next. Because this is a
large amount of information, it is organized into subordinate data blocks; in other
words, segment data blocks contain other kinds of data blocks.

A segment begins with a line containing the keyword \problemblock{*segment}, for
example,

\begin{lstlisting}[style=taosmanual]
*segment 23 2nd Stage Boost
\end{lstlisting}

The value following \problemblock{*segment} is the segment number, in this case
23. It is required and must be unique within the trajectory. An optional title, such as
\texttt{2nd Stage Boost}, follows the number.

The remainder of the segment definition consists of the data blocks listed in
\cref{tab:segment-data-blocks}. A segment block extends from its
\problemblock{*segment} keyword to the next trajectory- or problem-level keyword.
Usually the terminating keyword is another \problemblock{*segment} or a
\problemblock{*trajectory}. 
\Cref{tab:trajectory-data-blocks,tab:problem-data-blocks} list the trajectory- and
problem-level keywords.

\begin{table}[htbp]
  \centering
  \caption{Segment data blocks.}
  \label{tab:segment-data-blocks}
  \begin{tabularx}{0.88\linewidth}{@{}>{\ttfamily}lX@{}}
    \toprule
    \normalfont Data block & Description \\
    \midrule
    *aero      & Defines aerodynamic coefficients \\
    *constants & Provides values for user-defined variables \\
    *cg        & Defines the center of gravity \\
    *fly       & Defines a guidance rule \\
    *increment & Allows discontinuities in the vehicle state \\
    *inertial  & Aligns an inertial-platform coordinate system \\
    *integ     & Provides numerical-integration parameters \\
    *limits    & Defines limits on the flight path \\
    *prop      & Defines propulsive thrust and mass flow \\
    *rail      & Provides parameters for rail launches and sleds \\
    *reset     & Resets selected vehicle-state quantities \\
    *when      & Provides final conditions and what to do next \\
    \bottomrule
  \end{tabularx}
\end{table}

The following segment is an example of a rail-launched missile:

\begin{lstlisting}[style=taosmanual]
*segment 1 Rail Launch
  *integ dtprnt=0.10 dt=0.01
  *aero ca=(stage1-ca) cn=(stage1-cn)
  *prop thrust=(thrust1) mdot=(mdot1)
  *prop thrust=(t-recruits) mdot=(m-recruits)
  *rail launch cfstat=0.15 cfslid=0.01
  *when plength=25.0 goto 3
\end{lstlisting}

This segment contains an \problemblock{*integ} block giving the integration and
print step sizes, an \problemblock{*aero} block selecting aerodynamic tables, two
\problemblock{*prop} blocks selecting thrust and mass-flow tables, a
\problemblock{*rail} block defining rail-launch guidance and friction, and a
\problemblock{*when} block defining the final condition. The
\problemblock{*when} block also specifies how the trajectory continues; here,
segment 3 is computed next.

The following subsections describe the input parameters and format for each segment
data block. Additional complete segment examples are given in
\cref{sec:problem-examples}.
